
\section{Feature Overview}
	In this section an overview of the used features is given.
	The quantities used for the feature vector are depicted in Tab.~\ref{tab:Feature}
	and are generated with the GFN2--xTB method. The features are mostly atom centered quantities
	and are based on geometry, density or energy.
	The features were calculated for the systems in the APF CS.\\
	There are quantities, which also have an extended version.
	If this is the case, a tick is set at Extension.
	The quadrupole moment matrix $\bm{\Uptheta}$ was reduced to a vector by the multiplication
	with a one vector $\underline{1}$ (s.~Eq~\eqref{eq:redTheta}).
	\begin{align}
		\label{eq:redTheta}
		\bm{\uptheta} = \bm{\Uptheta} \times \underline{1}
	\end{align}

	\begin{table}[H]
		\centering
		\caption{List of quantities used for the ML models as features.
				The quantities are divided into geometry,
				density and energy based quantities.}
			\begin{tabular}{l l l}
				\hline
				Feature &  Description& Extension\\
				\hline
				\textbf{Geometry} & & \\
				$CN$ & Coordination Number & \checkmark \\
				\hline

				\textbf{Density} && \\
				$\bm{\upmu}$ & Dipole Moment & \checkmark \\
				$\bm{\upmu}_\mathrm{Z}$ & Dipole Moment, only Z direction & \\
				$\bm{\uptheta}$ & Quadrupole Moment & \checkmark\\
				\hline

				\textbf{Energy} && \\
				$E_\mathrm{resp}$ & Response& \\
				$E_\mathrm{gap}$ & HOMO -- LUMO gap& \checkmark\\

				$\mu_\mathrm{pot}$& Chemical Potential & \checkmark\\

				$\epsilon_\mathrm{HOAO}$ & Atomic HOMO & \checkmark \\

				$\epsilon_\mathrm{LUAO}$ & Atomic LUMO & \checkmark \\

				$E_\mathrm{rep}$ & Repulsion & \\
				$E_\mathrm{EHT}$ & extended Hückel theory& \\
				$E_\mathrm{disp}^{(2)}$ & Dispersion 2$^\mathrm{nd}$ order & \\
				$E_\mathrm{disp}^{(3)}$ & Dispersion 3$^\mathrm{rd}$ order & \\
				$E_\mathrm{IES+IXC}$ & Isotropic ES and XC & \\
				$E_\mathrm{AES}$ & Anisotropic ES& \\
				$E_\mathrm{AXC}$ & Anisotropic XC & \\
				$E_\mathrm{tot}$ & Total Energy & \\
				\hline 
			\end{tabular}
		\label{tab:Feature}
	\end{table}
	Each quantity is atom specific, thus each quantity is there from atom $A$ and $B$.
	The quantities shown in Tab.~\ref{tab:Feature} are all further processed for the use as features.
	From the quantities $x_A$ and $x_B$ the arithmetic mean,
	the product and absolute difference is calculated (s. Eq.~\eqref{eq:FeatureProcess}).
	The new quantities $X_\mathrm{Arith}$,
	$X_\mathrm{Prod}$ and $X_\mathrm{AbsDiff}$ are used as the features for the ML model.
	\begin{subequations}
		\label{eq:FeatureProcess}
		\begin{align}
			X_\mathrm{Arith}&= \dfrac{x_A + x_B}{2}\\
			X_\mathrm{Prod}&= x_A \cdot x_B \\
			X_\mathrm{AbsDiff}&= | x_A - x_B | 
		\end{align}
	\end{subequations}
	So far the feature $X$ is build up by $X = \{X_\mathrm{Arith},X_\mathrm{Prod},X_\mathrm{AbsDiff}\}$.
	If such feature is connected to the hessian matrices,
	it is independent for every element in a hessian matrix block $\textbf{H}_{AB}$.
	Hence, the ML model is not able to distinguish the elements.
	Two ways were developed and tested, which distinguish the elements.
	At first, the feature is increased by indices specific for each element. The indices are shown in Tab.~\ref{tab:Indices}
	\begin{table}[H]
		\centering
		\caption{Indices used for the different elements in a hessian matrix block to differ between the elements.}
			\begin{tabular}{c c c c | c }
				\hline
				& &  Indexed& & Permuted \\
				Hessian element &  $x^\mathrm{idx}$ & $y^\mathrm{idx}$ & $z^\mathrm{idx}$ & Index\\
				\hline
				\\[-1em]
				$\dfrac{\partial^2 E}{\partial x_A \partial x_A}$ & 2 & 0 & 0 & 0\\
				\\[-1em]
				$\dfrac{\partial^2 E}{\partial x_A \partial y_A}$ & 1 & -1 &  0 & 1\\
				\\[-1em] 
				$\dfrac{\partial^2 E}{\partial x_A \partial z_A}$ & 1 & 0 &  -1 & 2\\
				\\[-1em]
				$\dfrac{\partial^2 E}{\partial y_A \partial x_A}$ & -1 & 1 &  0 & 1\\
				\\[-1em]
				$\dfrac{\partial^2 E}{\partial y_A \partial y_A}$ & 0 & 2 &  0 & 0\\
				\\[-1em]
				$\dfrac{\partial^2 E}{\partial y_A \partial z_A}$ & 0 & 1 &  -1 & 2\\
				\\[-1em]
				$\dfrac{\partial^2 E}{\partial z_A \partial x_A}$ & -1 & 0 &  1 & 2\\
				\\[-1em]
				$\dfrac{\partial^2 E}{\partial z_A \partial y_A}$ & 0 & -1 &  1 & 2\\
				\\[-1em]
				$\dfrac{\partial^2 E}{\partial z_A \partial z_A}$ & 0 & 0 &  2 & 3\\
				\\[-1em]
				\hline 
			\end{tabular}
		\label{tab:Indices}
	\end{table}

	The feature vector is expanded with these indices
	$ X  = \{x^\mathrm{idx},y^\mathrm{idx},z^\mathrm{idx},X_\mathrm{Arith},X_\mathrm{Prod},X_\mathrm{AbsDiff}\}$.
	This labeling is called \textbf{indexed} (idx).\\
	In another approach, the hessian matrix is further rotated,
	so that the hessian elements are on the same spot for every prediction.
	In the same way, the feature quantities are adapted. The adapted matrix and coordinate system are shown in the Appendix 

	
	\begin{table}[H]
		\centering
		\caption[]{Permutation of the coordinates for each Hessian element.}
			\begin{tabular}{c c c c c}
				\hline
				Hessian Element & x & y & z \\
				\hline \\[-1em]
				$\dfrac{\partial^2 E}{\partial x_A \partial x_A}$ & x & y & z & \\ [+0.7em]
				$\dfrac{\partial^2 E}{\partial x_A \partial y_A}$ & y & x & $-$z & \\[+0.7em]
				$\dfrac{\partial^2 E}{\partial x_A \partial z_A}$ & z & x & y & \\[+0.7em]

				$\dfrac{\partial^2 E}{\partial y_A \partial x_A}$ & x & y & z & \\[+0.7em]
				$\dfrac{\partial^2 E}{\partial y_A \partial y_A}$ & y & z & x & \\[+0.7em]
				$\dfrac{\partial^2 E}{\partial y_A \partial z_A}$ & z & y & $-$x & \\[+0.7em]

				$\dfrac{\partial^2 E}{\partial z_A \partial x_A}$ & x & z & $-$y & \\[+0.7em]
				$\dfrac{\partial^2 E}{\partial z_A \partial y_A}$ & y & z & x & \\[+0.7em]
				$\dfrac{\partial^2 E}{\partial z_A \partial z_A}$ & z & x & y & \\[+0.7em]
				\hline

			\end{tabular}
		\label{tab:Permutation}
	\end{table}



	After that, an index is put onto the features again. 
	In this case though, the hessian matrix elements with a derivative in respect to the $z$ coordinate get the same index.
	Furthermore, the diagonal and off--diagonal elements get different indices.
	This labeling is called \textbf{permuted} (perm).
	%CCCCCCCCCCCCCCCCCCCCCCCCCCCCCCCCCCCCCCCCCCCCCCCCCCCCCCCCCCCCCCCCCCCCCCCCCCCCCCCCCCCCCCCCCCCCC
	\section{Transformation of Coordinates}
	The hessian matrix and some features depend on the orientation of the coordinate~system~(CS).
	For the ML model a consistent orientation for every atom pair is necessary.
	Furthermore, the hessian elements shall be learned atom pair wise.
	Thus, the orientation of the molecule is brought into atom pair focused CS,
	where the origin of the CS is in the center of the considered atom pair,
	while both atoms are on the z--Axis.
	In this section,
	the transformation of the initial coordinate system into the atom pair focused CS will be discussed.\\
	The initial CS is not fixed to any constraints.
	Therefore, the initial CS is transformed into the main inertia system,
	from which the transformation to the atom pair focused coordinate systems is performed.\\
	At first a translation of the origin of the initial CS to the center of mass is done by Eq.~\eqref{eq:TransIinitCenterMass}.
	
	\begin{align}
		\textbf{r}^\mathrm{s}_A =  \textbf{r}^\mathrm{init}_A - \textbf{s}^\mathrm{init}
		\label{eq:TransIinitCenterMass}
	\end{align}
	Where $\textbf{r}^\mathrm{s}_A$ is the coordinate vector of an atom $A$ in the mass centered CS, $\textbf{r}^\mathrm{init}_A$ the coordinate vector in the initial CS and $\textbf{s}^\mathrm{init}_A$ is the center of mass in the initial CS, and is defined in Eq.~\eqref{eq:CenterOfMass}.
	
	\begin{align}
		\textbf{s}^\mathrm{init} = \dfrac{1}{M} \sum_{A}^{N} m_A \cdot \textbf{r}^\mathrm{init}_A
		\label{eq:CenterOfMass}
	\end{align}
	Here $m_A$ is the mass of atom $A$, $M$ is the mass of the molecule and $N$ the number of atoms in the molecule.
	Once the origin of the new CS is determined,
	the orientation of its axes is arranged to the moment of inertia as computed by Eq.~\eqref{eq:MomentOfInertia}.
	
	\begin{align}
		\textbf{I}^\mathrm{s} = \left(\begin{matrix}
						I_{xx} & I_{xy} & I_{xz} \\
						I_{yx} & I_{yy} & I_{yz} \\
						I_{zx} & I_{zy} & I_{zz} \\ 
				   	\end{matrix} \right) 
			 ; I_{ab} = \begin{cases}
				\sum_{A}^{N} m_A ( b_A^2 + c_A^2 ) ~~\mathrm{if}~ a = b\\
			   	 - \sum_{A}^{N} m_A \cdot a_A \cdot b_A  ~\mathrm{if}~ a \neq b \\
			   	\end{cases} a,b \in x,y,z
		\label{eq:MomentOfInertia}
	\end{align}
	
	For the computation of the elements $I_{ab}$ of the inertia tensor $\textbf{I}^\mathrm{s}$ the diagonal
	elements and non--diagonal elements are divided in two cases.
	Here, $a$ and $b$ are variables for the spatial directions of the coordinates~$x,y$ and $z$.
	For the rotation of the mass centered origin CS into the main inertia system,
	the matrix $\textbf{R}_\mathrm{mi}^\mathrm{init}$ of eigenvectors of the inertia tensor $\textbf{I}^\mathrm{s}$ is needed.
	For that, the eigenvectors $\textbf{R}_\mathrm{mi}^\mathrm{init}$ are computed
	by solving the eigenvalue equation.
	The eigenvalues are then used to get the eigenvectors.
	To assure, that the matrix $\textbf{R}_\mathrm{mi}^\mathrm{init}$ is a right--handed matrix the determinant is calculated.
	If the determinant is negative, the sign of the third eigenvector is changed. Furthermore, it needs to be assured,
	that the eigenvector is consistently rotated in the same direction, independent of the outgoing position.
	The chosen constraint is, that the highest value of the eigenvectors are positive,
	if this is not the case the sign of the eigenvector and the next eigenvector in $\textbf{R}_\mathrm{mi}^\mathrm{init}$ is changed.\\
	The rotation of the coordinates is done by Eq.~\eqref{eq:rotCoord}.
	\begin{align}
		\textbf{r}_A^\mathrm{mi} = \textbf{R}_\mathrm{mi}^\mathrm{init} \textbf{r}_A^\mathrm{s}
		\label{eq:rotCoord}
	\end{align}
	Here, $\textbf{r}_A^\mathrm{mi}$ is the coordinate of an atom $A$ in the main inertia CS.
	The rotation is done for every atom in the molecule. After that operation,
	the molecule is in the main inertia CS.\\ The hessian matrix also depends
	on the CS and therefore needs to be rotated, too.
	A translation is not needed. Eq.~\eqref{eq:rotHessian} describes
	the rotation of the hessian matrix in the initial CS to the main inertia system.
	For this a $3N \times 3N$ rotation matrix $\textbf{P}_\mathrm{mi}^\mathrm{init}$ needs to be constructed,
	where the diagonal 3 x 3 blocks are $\textbf{R}_\mathrm{mi}^\mathrm{init}$.
	This matrix is shown in Eq.~\eqref{eq:rotMatrixHessian}.
	\begin{align}
		\label{eq:rotMatrixHessian}
		\textbf{P}_\mathrm{mi}^\mathrm{init} = \left(\begin{matrix}
							\textbf{R}_\mathrm{mi}^\mathrm{init} &\text{\Large\textbf{0}}& && \text{\Large\textbf{0}}  \\
							\text{\Large\textbf{0}}& \textbf{R}_\mathrm{mi}^\mathrm{init} & & &\text{\Large\textbf{0}} \\
							\text{\Large\textbf{0}} & & \ddot& &\\
							 & & &\ddot & \text{\Large\textbf{0}}\\
							\text{\Large\textbf{0}}& & & \text{\Large\textbf{0}} &\textbf{R}_\mathrm{mi}^\mathrm{init}
							\end{matrix}\right)
	\end{align}

	\begin{align}
		\textbf{H}^\mathrm{mi} =\textbf{P}_\mathrm{mi}^\mathrm{init} \times \textbf{H}^\mathrm{init} \times {\left(\textbf{P}_\mathrm{mi}^\mathrm{init}\right)}^\intercal
		\label{eq:rotHessian}
	\end{align}
	The hessian matrix in the initial system is $\textbf{H}^\mathrm{init}$,
	while the hessian matrix in the main inertia system is $\textbf{H}^\mathrm{mi}$.\\
	In the following, the transformation of the main inertia CS to the atom pair focused CS will be discussed.
	The transformation starts with the translation of the origin to the center of an atom pair.
	The translation is described with Eq.~\eqref{eq:TranslationInertiaAPF}.\\
	\begin{align}
		\textbf{r}^\mathrm{mi,c}_C = \textbf{r}^\mathrm{mi}_C -
		 \left(\dfrac{1}{2} \textbf{r}^\mathrm{mi}_A + 
		 \dfrac{1}{2}\textbf{r}^\mathrm{mi}_B \right) 
		 \mathrm{for}~A,B,C \in A, B, ..., N
		\label{eq:TranslationInertiaAPF}
	\end{align}
	After this operation, the origin of the CS is in the center of the atom pair.
	To fix the orientation, the atom pair is aligned along the z--axis.
	The Euler angles are used to determine the rotation matrix.
	The rotation of the CS depends on the considered atom pair.
	It is divided into whether $A = B$ or $A \neq B$.
	At first the construction of the rotation matrix for the case $A \neq B$ will be discussed.\\
	The z--axis of the atom pair focused CS is aligned to $\textbf{r}^\mathrm{mi,c}_A$.
	The x--axis is defined in Eq.~\eqref{eq:x_axis}, where $\textbf{z}$ is a unit vector in z--direction of the atom pair focused CS,
	and $\bm{\upmu}_A$ is the dipole moment of the atom $A$.
	\begin{align}
		\textbf{x} = \textbf{z} \times ( \mathrm{\bm{\upmu}_{A}} + \mathrm{\bm{\upmu}_{B}})
		\label{eq:x_axis}
	\end{align}
	The cutting line $\textbf{L}$ of the xy--plane of the atom pair focused CS and xy--plane of the centered main inertia CS
	is described with Eq.~\eqref{eq:EulerLL}.
	\begin{align}
		\label{eq:EulerLL}
		\mathbf{L} = \mathbf{x} \times \mathbf{z}^\mathrm{mi,c}
	\end{align}
	The cutting line $\textbf{L}$, the axes $\textbf{x},\textbf{z},\textbf{x}^\mathrm{mi,c}$ and $ \textbf{z}^\mathrm{mi,c}$ are then used for the computation of the Euler angles  $\alpha, \beta$ and $\gamma$ via Eq.~\eqref{eq:EulerAnglesDiscrete}.
	\begin{subequations}
	\label{eq:EulerAnglesDiscrete}
	\begin{align}
		\alpha^{0}  &= \arccos{ \left( \dfrac{\textbf{L} \cdot \textbf{x}^\mathrm{mi,c}}{|\textbf{L}|\cdot |\textbf{x}^\mathrm{mi,c}|}\right) }\\
		\beta &= \arccos{ \left( \dfrac{\textbf{x}	\cdot \textbf{z}^\mathrm{mi,c}}{|\textbf{x}|\cdot |\textbf{z}^\mathrm{mi,c}|} \right)} \\
		\gamma^{0} &= \arccos{ \left( \dfrac{\textbf{L} \cdot  \textbf{x}}{|\textbf{L}|\cdot |\textbf{x}|}\right) }
	\end{align}
\end{subequations}
	For the right rotation, it has to be considered,
	whether the vectors $\textbf{x}^\mathrm{mi,c}, \textbf{z}^\mathrm{mi,c}$ and $\textbf{L}$
	span a right--handed coordinate system or left--handed.
	This is done by constructing a matrix of these vectors and computing the determinant.
	The determinant can be used, as it is equal to the triple product for a $3\times 3$ matrix.
	
	\begin{align}
		\textbf{A} &= (\textbf{x},\textbf{z},\textbf{L}) \\
		\textbf{G} &= (\textbf{L},\textbf{x}^\mathrm{mi,c},\textbf{z}^\mathrm{mi,c})
	\end{align}

\begin{align}
	\alpha &= \begin{cases}
		 \alpha^{0}  ~~~~~~~~\mathrm{if}~ \det (\textbf{A}) = 1  \\
	2\pi - \alpha^{0} ~\mathrm{if}~ \det (\textbf{A}) = -1 \\
	\end{cases} \\
	\gamma   &= \begin{cases}
		\gamma^{0} ~~~~~~~~\mathrm{if}~ \det (\textbf{G}) = -1  \\
		2\pi - \gamma^{0} ~\mathrm{if}~ \det (\textbf{G}) = 1 \\
	\end{cases}
\end{align}
Via these constraints, the Euler angles $\alpha$ and $\gamma$ are chosen.
The Euler angles are then used to build up the rotation matrices via Eq.~\eqref{eq:rotMatZ} and  \eqref{eq:rotMatX}.\\
The complete rotation matrix is then constructed by the cross product of the single rotation matrices
and is shown in Eq.~\eqref{eq:RotMatFull}.
\begin{align}
	\label{eq:RotMatFull}
	\textbf{R}^0 = \textbf{R}_z(\gamma)  \times \textbf{R}_x(\beta)\times \textbf{R}_z(\alpha)
\end{align}

If the x--values of the rotated dipole moment $\bm{\upmu}$ are negative,
the matrix is rotated by 180$\degree$ around the z--axis. This is depicted in Eq.~\eqref{eq:rot180}.
\begin{align}
	\label{eq:rot180}
	\textbf{R}_\mathrm{mi}^\mathrm{APF} = \begin{cases}
		\textbf{R}^0 ~~~~~~~~~~~~~\mathrm{if}~ \bm{\upmu}^x > 0  \\
		 \textbf{R}^0 \times \textbf{R}_z(\pi) ~\mathrm{if}~ \bm{\upmu}^x < 0  \\
	\end{cases}
\end{align}
Lastly, the rotation of the translated coordinate system is done via Eq.~\eqref{eq:rotCS}.
\begin{align}
	\label{eq:rotCS}
	\textbf{r}^\mathrm{APF}_A = \textbf{R}_\mathrm{mi}^\mathrm{APF} \times \textbf{r}^\mathrm{mi,c}_A
\end{align} 
The coordinates $\textbf{r}^\mathrm{APF}_A$ are the final coordinates in the APF CS for the case of a heteronuclear atom pair.\\ 
For the case of a homonuclear atom pair,
the procedure is the same up to the translation of the coordinates into the centered inertial main CS.
The z--axes of APF CS and centered inertial main CS are aligned.
Therefore, the Euler angle $\beta$ is zero. Furthermore, the xy--planes are the same,
and there is no cutting line of the xy--planes.
Thus, only one angle $\alpha$ or $\gamma$ is necessary for the rotation.
Therefore, $\gamma$ is the angle between the x--axis of the centered main inertial CS
and the APF CS. To provide consistency in the rotation of the coordinates,
the x--axis is chosen in dependency of the dipole moment $\bm{\upmu}$ of the considered atom.
Eq.~\eqref{eq:xHomonuclear} is show how the x--axis is computed.
\begin{align}
	\label{eq:xHomonuclear}
	\textbf{x} = \textbf{z} \times \bm{\upmu}_A
\end{align}
This results in the computation of the angle $\gamma^0$ by Eq.~\eqref{eq:GammaHomonuclear}.

\begin{align}
	\label{eq:GammaHomonuclear}
		\gamma^{0} &= \arccos{ \left( \dfrac{\textbf{x}^\mathrm{mi,c} \cdot  \textbf{x}}{|\textbf{x}^\mathrm{mi,c}|\cdot |\textbf{x}|}\right) }
\end{align}
Between two vectors are two angles, which sum up to $2\pi$.
The rotation of the CS depends on the arrangement to another.
The angle is chosen by the following condition, show in Eq. \eqref{eq:GammaHomonuclearAdj}.
\begin{align}
	\label{eq:GammaHomonuclearAdj}
	\gamma = \begin{cases}
		\gamma^0 ~~~~~~~~\mathrm{if}~ \textbf{x}_y < 0\\
		2\pi - \gamma^0  ~\mathrm{if}~ \textbf{x}_y \geq 0 \\
		   \end{cases} 
\end{align}
The further procedure of the transformation is done in the same way as for the heteronuclear atomic pair.
The rotation matrix is computed via Eq.~\eqref{eq:RotMatFull}.
This followed by an adjustment of the rotation matrix,
shown in Eq.~\eqref{eq:rot180}.
The final rotation matrix $\textbf{R}_\mathrm{mi}^\mathrm{APF}$ is then
used on the coordinates $ \textbf{r}^\mathrm{mi,c}_A$ via Eq.~\eqref{eq:rotCS}.\\
In the APF CS the computation of the hessian was done. For each APF CS one $3\times 3$ hessian matrix block is of interest.
To rebuild the hessian matrix in the mi system, each hessian matrix block needs to undergo a rotation back to the mi system.
For this the all rotation matrices $\textbf{R}_\mathrm{mi}^\mathrm{APF}$ need to be stored.
The rotation back into the mi system is done by Eq.~\eqref{eq:rotHessBlock}.
\begin{align}
	\label{eq:rotHessBlock}
	\textbf{H}^\mathrm{mi}_\mathrm{AB} = {\left(\textbf{R}_\mathrm{mi}^\mathrm{APF}\right)}^\intercal
	 \times \textbf{H}^\mathrm{APF}_\mathrm{AB} \times  \textbf{R}_\mathrm{mi}^\mathrm{APF}
\end{align}
In the mi system, the hessian matrix $\textbf{H}^\mathrm{mi}$ is filled with the blocks $\textbf{H}^\mathrm{mi}_\mathrm{AB}$.
Before some quantities for measuring can be computed, the $\textbf{H}^\mathrm{mi}$ is projected and mass--weighted.
Thus, the $\textbf{H}^\mathrm{proj,w|mi}$ is build. 
The procedure of the projection and mass--weighting is described in the chapter Theory.

%  
%$
% The rotation of the hessian from the main inertial CS to the APF CS is done via Eq.~\eqref{eq:rotHessian}.
% The rotation matrix $\textbf{P}_\mathrm{mi}^\mathrm{APF}$
% is constructed the same as the rotation matrix $\textbf{P}_\mathrm{init}^\mathrm{mi}$.

\section{Quantities for Measuring}
The mean squared error~MSE was chosen for the quantification of the learning process of the ML models.
The ML models are trained with GFN2--xTB based information. Therefore, the reference is also GFN2--xTB.
In Eq.~\eqref{eq:MSE} is the MSE shown.
Here, $y$ corresponds to the hessian elements, calculated with GFN2--xTB and predicted with the ML model.
The number of data points is $n$.
\begin{align}
	\label{eq:MSE}
	\mathrm{MSE} =  \dfrac{1}{n} \sum_{i=1}^n \left( y_{i\mathrm{,ML}} - y_{i\mathrm{,xTB}} \right)^2
\end{align}
For the comparison of the different models (e.g. ETR--Case--Perm, RFR--Gen--Idx, etc.) the comparison is done by assuming a normal distribution.

For this the difference between the GFN2--xTB computation and ML model prediction is calculated.
Furthermore, the standard deviation and mean is computed. After that a normal distribution is assumed.
This normal distribution is computed with Eq.~\eqref{eq:NormalDistribution}.
\begin{align}
	\label{eq:NormalDistribution}
	f(x)  = \dfrac{1}{\sigma \sqrt{2\pi}} \exp\left({-\dfrac{1}{2} \left(\dfrac{x-\mu}{\sigma}\right)^2}\right)
\end{align}

The physical quantities, that were investigated are the elements of the hessian matrix $\textbf{H}^\mathrm{mi}$,
the wave number $\tilde{\nu}$, the force constant $k$ and the zero point vibrational energy ZPVE.
The wave numbers are calculated by Eq.~\eqref{eq:WN},
the force constants by Eq.~\eqref{eq:FC} and the ZPVE by Eq.~\eqref{eq:ZPVE}.

Here, $\lambda$ are the eigenvalues of the mass weighted projected hessian matrix $\textbf{H}^\mathrm{proj,w|mi}$ and $m_i$ the mass of each atom.
For the transformation into SI--units, the hartree Energy $E_\mathrm{h}$,
planck constant $h$, reduced planck constant $\hbar$, speed of light $c$ and the atomic mass unit $u$ are required.

\begin{subequations}
\begin{align}
	\label{eq:WN}
	\tilde{\nu} &= \sqrt{|\lambda|} \cdot \dfrac{E_\mathrm{h}}{\hbar \cdot c}\\
	\label{eq:FC}
	k &=   \dfrac{E_\mathrm{h}^2}{\hbar^2 \cdot u} \cdot \sum^M_i \dfrac{1}{m_i} \lambda\\
	\label{eq:ZPVE}
	\mathrm{ZPVE} &= \dfrac{E_\mathrm{h}}{h} \sum_i^n \sqrt{|\lambda_i|} 
\end{align}
\end{subequations}



\section{Computational Methods}

All electronic structure computations were performed with GFN2--xTB\supercite{Bannwarth2019_GFN2xTB}, for this \textit{xtb}\supercite{Bannwarth2021_xTB} was used.
The computation of the information for the feature vector and target vector are done with GFN2--xTB level calculation.
For this the GFN2--xTB code was used. For transformation of the coordinates and the ML models python was used.
Especially, the packages \textit{sklearn}\supercite{scikit-learn}, \textit{numpy}\supercite{numpy} and \textit{pandas}\supercite{pandas.Paper,pandas.software} were used for this project.
